\chapter{Surface tension and gravity}\label{ch:forces}
Thesis Chapters~4 and~7 discuss surface tension, the Young--Laplace jump,
and the mimetic gravity model~\cite{remmerswaal2023}. This chapter will map
those ideas to the available balanced forces and curvature handling.

\section{CSF and conservative ghost-fluid forces}
The coefficient $\sigma$ and curvature $\kappa$ determine the pressure jump;
face curvature weighting and its fallback will be specified here.
\begin{deviation}{Aperture surface tension}
The thesis's aperture surface-tension model is not implemented; the supported choices are CSF and CGFM.
\end{deviation}
\begin{deviation}{Geometric-distance GFM}
The thesis's nonconservative geometric-distance GFM is not a supported surface-tension choice.
\end{deviation}
\begin{deviation}{Force stage}
Surface tension is added to the predictor velocity, whereas the thesis writes it inside the scaled pressure-gradient operator.
\end{deviation}
\begin{implnote}
\apitype{SurfaceTension}, \apitype{CSF}, \apitype{CGFM}, \apifun{surface_tension!}.
\end{implnote}

\section{Gravity and well-balancedness}
The gravity vector $\+g$ acts through an interface-aware potential. The
well-balanced relation to the thesis gravity model remains to be checked.
\begin{implnote}
\apitype{TwoPhase}, \apifun{gravity!}.
\end{implnote}
